%=============================================================================
\section{Introduction}
\label{sec:intro}
%=============================================================================

Cross-chain bridges move assets between heterogeneous ledgers and have become core infrastructure of decentralized finance (DeFi)~\cite{robinson2021survey,ou2022overview}.
They are also a preferred laundering route.
Bridge exploits account for many of the largest thefts in the ecosystem~\cite{lee2023sok}, and Elliptic estimates that more than \$21 billion of criminal and high-risk funds had passed through DEXs, bridges and swap services by 2025, a threefold rise over two years~\cite{elliptic2025crosschain}.
After an exploit, the attacker typically moves the proceeds within hours: fresh intermediary wallets, token swaps into liquid assets, fan-out splits, and finally a bridge deposit, a mixer or an exchange, after which the trail is lost or frozen funds can no longer be recovered.
MITRE's AADAPT framework names these steps cross-chain hopping, layering, peel chains and anonymizing services~\cite{mitre2025aadapt}.

Research has made steady progress on \emph{reconstructing} such flows.
CONNECTOR and ABCTracer associate deposits and withdrawals across bridges~\cite{lin2025connector,lin2025abctracer}, BridgeGuard and XChainWatcher detect attack transactions and rule violations on bridges~\cite{wu2025safeguarding,augusto2025xchainwatcher}, and AMLGuard combines semantic transaction analysis with retrieval-augmented language-model reasoning to trace laundering paths on 82 real cases~\cite{wu2026amlguard}.
These systems answer a forensic question, \emph{where did the money go?}, and they work on complete or nearly complete histories.
A compliance team, a bridge operator or an exchange facing an incident \emph{while it unfolds} needs a different answer: given only the actions observed so far, is this flow on a laundering path, and how much time is left before it reaches an exit that cannot be undone?

We call this problem \emph{early detection of cross-chain laundering} and study it on real incidents.
It differs from forensic tracing in three ways.
First, the detector sees only a prefix of the trajectory and must score it without any information that depends on how the trajectory ends.
Second, most prefixes of benign activity look like laundering: a user who bridges, swaps and splits is not a thief, so the detector must be evaluated against benign trajectories that pass through the same bridges at the same time.
Third, a score is only useful if it is calibrated and explained, because a human decides whether to freeze or escalate.

We present \textbf{CAGI-ED} (Cross-chain Attack Graph Intelligence for Early Detection), a CPU-only pipeline that monitors a seed address, decodes every new action into a DeFi semantic type (transfer, swap, bridge, lending, split, merge, mixer or exit), expands the tainted value flow, and re-scores the growing trajectory after each action.
The scorer uses \emph{typed temporal motifs}, features that describe what kind of actions happened, in which order and how fast, and turns them into a calibrated probability with an explanation in AADAPT terms.

\paragraph{Contributions.}
\begin{itemize}
\item \textbf{Problem and protocol.} We formalize prefix-based early detection with an explicit endpoint, lead time and per-trajectory false-alert rate, and we separate online horizons (available to a live monitor) from relative checkpoints (a retrospective diagnostic). We evaluate under leave-one-incident-out with incident-level bootstrap intervals, and we report a leak that we found and removed in our own earlier pipeline (Sect.~\ref{sec:problem},~\ref{sec:setup}).
\item \textbf{System.} An end-to-end architecture from multi-chain explorer APIs to explained alerts: an evidence cache with checksums, a semantic decoder over a verified protocol map, a bounded value-flow trajectory builder, 37 prefix-safe features in three groups, and a calibrated gradient-boosted scorer with TreeSHAP attributions mapped to MITRE AADAPT (Sect.~\ref{sec:method}).
\item \textbf{Dataset.} 15 publicly documented cross-chain incidents (2021--2025) on Ethereum, BNB Smart Chain and Arbitrum, each traced from its seed, with 393 hard-negative trajectories mined from the same bridges, DEXs and mixers in the same time window. We release the decoded trajectories, the raw evidence cache and all scripts (Sect.~\ref{sec:data}).
\item \textbf{Evidence.} A controlled study that holds the learner fixed and varies only the representation, results at online horizons from one minute to one day, lead time and false alerts per threshold, ablations, calibration, a mimicry stress test and cost on commodity hardware (Sect.~\ref{sec:results}).
\end{itemize}
